\section{Proposed Defense Methods}
\label{sec:method}

\subsection{Control-Plane Attack Insight}

Before describing our defenses, we motivate the frequency-domain approach by
establishing the nature of adversarial attacks on AMC. We conduct a two-track
experiment to separate \emph{classification failure} from \emph{waveform
corruption}. In Track~A, we apply CW and EAD attacks to real RML2016.10a
signals and record classification accuracy. In Track~B, we generate synthetic
bursts with known payload bits and CRC-8 checksums using a full TX/RX chain
(pulse shaping, AWGN channel, matched filtering, pilot-based equalization),
transfer the adversarial perturbation vectors from Track~A, and evaluate CRC
pass rate under four demodulation conditions:

\begin{enumerate}
  \item Clean+Oracle: clean signal with true modulation known;
  \item Clean+AMC: clean signal, modulation predicted by AWN;
  \item Adv+Oracle: adversarial signal with true modulation known;
  \item Adv+AMC: adversarial signal, modulation predicted by AWN.
\end{enumerate}

Table~\ref{tab:crc} shows CRC pass rates at SNR~18~dB under CW attack for
eight digital modulations. The Adv+Oracle column reveals that the underlying
signal remains largely intact under adversarial perturbation: CRC pass rates
stay above 71\% when the correct demodulator is applied. The catastrophic
failure manifests only in the Adv+AMC column (0--39\%), where the classifier
selects the wrong demodulator.

\begin{table}[t]
\caption{CRC Pass Rate (\%) at SNR 18~dB --- CW Attack. Adversarial
  perturbation drops AMC accuracy to near 0\% (Adv+AMC) but CRC integrity
  under oracle demodulation (Adv+Oracle) remains above 71\%, confirming
  that attacks target the classification decision, not the physical signal.}
\label{tab:crc}
\centering
\small
\begin{tabular}{lcccc}
\toprule
Mod & Clean+Orc & Clean+AMC & Adv+Orc & Adv+AMC \\
\midrule
BPSK  & 100.0 & 97.0 & 100.0 & 39.0 \\
QPSK  & 100.0 & 98.5 & 100.0 &  1.5 \\
8PSK  & 100.0 & 98.0 & 100.0 &  0.5 \\
QAM16 & 100.0 & 94.5 &  99.5 &  0.5 \\
QAM64 &  96.0 & 94.0 &  76.0 &  0.5 \\
PAM4  & 100.0 & 98.5 &  71.0 &  2.0 \\
CPFSK & 100.0 & 100.0 & 100.0 & 18.5 \\
GFSK  &  99.5 & 99.0 &  99.5 &  0.0 \\
\bottomrule
\end{tabular}
\end{table}

This \emph{control-plane attack} insight motivates our frequency-domain
approach: adversarial perturbations are spectral artifacts that can in principle
be attenuated by frequency-domain filtering without destroying the underlying
modulation waveform.

\subsection{FFT Top-\texorpdfstring{$K$}{K} Filtering}

The fundamental building block of our defense is FFT Top-$K$ filtering,
which retains the $K$ largest-magnitude spectral components and zeros all
others. Formally, for a single IQ sample $\mathbf{x} \in \mathbb{R}^{2 \times T}$:

\begin{enumerate}
  \item \textbf{Normalize:} $\tilde{\mathbf{x}} = (\mathbf{x} + x_{\mathrm{off}})
    / x_{\mathrm{scale}}$, mapping amplitudes to $[0, 1]$.
  \item \textbf{FFT:} $\mathbf{X} = \text{rFFT}(\tilde{\mathbf{x}})$, producing
    a complex spectrum $\mathbf{X} \in \mathbb{C}^{2 \times (T/2+1)}$.
  \item \textbf{Top-$K$ mask:} For each channel $c \in \{I, Q\}$, compute
    magnitudes $|\mathbf{X}_c|$ and construct binary mask $\mathbf{m}_c$
    that retains the $K$ largest entries and zeros the rest.
  \item \textbf{Inverse FFT:} $\hat{\tilde{\mathbf{x}}} =
    \text{irFFT}(\mathbf{X} \odot \mathbf{m})$.
  \item \textbf{Denormalize:} $\hat{\mathbf{x}} = \hat{\tilde{\mathbf{x}}}
    \cdot x_{\mathrm{scale}} - x_{\mathrm{off}}$.
\end{enumerate}

In compact notation:
\begin{equation}
  \hat{\mathbf{x}} = \text{irFFT}\!\bigl(\mathrm{TopK}\bigl(\text{rFFT}(\tilde{\mathbf{x}}),\, K\bigr)\bigr),
  \label{eq:topk}
\end{equation}
where the normalization and denormalization are implicit. The rationale is that
adversarial perturbations, which are added in the amplitude domain, tend to
distribute energy across the full spectrum; retaining only the dominant $K$
components removes much of this diffuse adversarial energy while preserving the
modulation-carrying low-frequency components.

\subsection{Spectral-Gated Defense}
\label{sec:sg}

FFT Top-$K$ filtering is effective for narrowband modulations (BPSK, QPSK,
QAM, etc.) but completely fails for wideband modulations such as AM-SSB,
whose spectral energy is spread uniformly across frequencies. Top-$K$ filtering
of AM-SSB destroys the signal itself, reducing recovery accuracy to 0\% even
when $K$ is generous.

We exploit a robust spectral feature, \emph{spectral flatness}, to detect
wideband signals without a separate classifier. Spectral flatness is defined as
the ratio of the geometric mean to the arithmetic mean of the power spectrum:
\begin{equation}
  \text{SF}(\mathbf{x}) = \frac{\exp\!\bigl(\tfrac{1}{N}\sum_{k}\ln |X_k|^2\bigr)}
                               {\tfrac{1}{N}\sum_{k}|X_k|^2},
  \label{eq:flatness}
\end{equation}
where $X_k$ are the rFFT coefficients and $N = T/2 + 1$ is the number of
positive-frequency bins. For narrowband modulations at SNR~18~dB, spectral
flatness is below 0.04 (all energy concentrated in few bins). For AM-SSB,
flatness is approximately 0.55 (energy spread uniformly across frequencies).
This gap of more than $18\times$ provides a robust routing threshold with
negligible misclassification between the two routes.

Algorithm~\ref{alg:gated} presents the spectral-gated defense. The key
efficiency insight is that the FFT computation is \emph{shared}: the same
FFT coefficients are used for both the flatness check and the Top-$K$
operation, requiring no additional computation.

\begin{algorithm}[t]
\caption{Spectral-Gated Defense}
\label{alg:gated}
\begin{algorithmic}[1]
\Require Input $\mathbf{x} \in \mathbb{R}^{N \times 2 \times T}$,
  Top-$K$ parameter $K$, quantization levels $L$, flatness threshold $\tau$
\Ensure Defended signal $\hat{\mathbf{x}}$
\State Normalize: $\tilde{\mathbf{x}} \leftarrow (\mathbf{x} + x_{\mathrm{off}}) / x_{\mathrm{scale}}$
\State $\mathbf{X} \leftarrow \text{rFFT}(\tilde{\mathbf{x}})$
  \hfill \textit{// shared FFT for all samples}
\For{each sample $i = 1, \ldots, N$}
  \State $\text{SF}_i \leftarrow \text{SpectralFlatness}(\mathbf{X}_i)$
    \hfill \textit{// Eq.~\eqref{eq:flatness}}
  \If{$\text{SF}_i > \tau$}
    \hfill \textit{// wideband (AM-SSB)}
    \State $\hat{\tilde{\mathbf{x}}}_i \leftarrow \text{Quantize}(\tilde{\mathbf{x}}_i, L)$
  \Else
    \hfill \textit{// narrowband (all others)}
    \State $\hat{\tilde{\mathbf{x}}}_i \leftarrow \text{irFFT}(\mathrm{TopK}(\mathbf{X}_i, K))$
      \hfill \textit{// reuse FFT}
  \EndIf
\EndFor
\State Denormalize: $\hat{\mathbf{x}} \leftarrow \hat{\tilde{\mathbf{x}}} \cdot x_{\mathrm{scale}} - x_{\mathrm{off}}$
\State \Return $\hat{\mathbf{x}}$
\end{algorithmic}
\end{algorithm}

Default parameters are $K = 20$, $L = 32$, and $\tau = 0.4$.
The quantization routing recovers AM-SSB from 0\% to 20--65\% accuracy under
CW and EAD attacks, a modulation where Top-$K$ alone achieves 0\% recovery.

\subsection{Adaptive-\texorpdfstring{$K$}{K} Defense}
\label{sec:adaptivek}

Fixed-$K$ filtering applies the same spectral truncation to all signals
regardless of their inherent bandwidth, which can over-filter wideband
digital modulations or under-filter narrowband signals under low-SNR attacks.
The Adaptive-$K$ defense selects the number of retained spectral components
per sample based on the spectral magnitude distribution.

Algorithm~\ref{alg:adaptivek} describes the per-sample $K$ selection procedure.
The key idea is that the spectral magnitude profile of a clean modulated signal
has a characteristic shape: a few dominant components carry most of the energy,
followed by a rapid decay toward noise floor. The \emph{knee point} of this
profile---where the cumulative energy first exceeds a threshold $\eta$---defines
the natural $K$ for that signal. Adversarial perturbations, by distributing
energy across many bins, push the knee point to larger $K$ values; the adaptive
method corrects for this by anchoring $K$ to the cumulative energy threshold
rather than a fixed count.

\begin{algorithm}[t]
\caption{Adaptive-$K$ Spectral Defense}
\label{alg:adaptivek}
\begin{algorithmic}[1]
\Require Input $\mathbf{x} \in \mathbb{R}^{N \times 2 \times T}$,
  energy threshold $\eta \in (0, 1)$, $K_{\min}$, $K_{\max}$
\Ensure Defended signal $\hat{\mathbf{x}}$
\State Normalize: $\tilde{\mathbf{x}} \leftarrow (\mathbf{x} + x_{\mathrm{off}}) / x_{\mathrm{scale}}$
\State $\mathbf{X} \leftarrow \text{rFFT}(\tilde{\mathbf{x}})$
\For{each sample $i = 1, \ldots, N$}
  \State $\mathbf{m}_i \leftarrow |\mathbf{X}_i|$ \hfill \textit{// per-channel magnitudes, flattened}
  \State Sort $\mathbf{m}_i$ descending: $\mathbf{m}_i^{\downarrow}$
  \State $E_{\mathrm{total}} \leftarrow \sum_k (\mathbf{m}_i^{\downarrow}[k])^2$
  \State $K_i \leftarrow \min\{k : \sum_{j=1}^{k}(\mathbf{m}_i^{\downarrow}[j])^2
    \geq \eta \cdot E_{\mathrm{total}}\}$
  \State $K_i \leftarrow \mathrm{clip}(K_i,\; K_{\min},\; K_{\max})$
  \State $\hat{\tilde{\mathbf{x}}}_i \leftarrow \text{irFFT}(\mathrm{TopK}(\mathbf{X}_i, K_i))$
\EndFor
\State Denormalize: $\hat{\mathbf{x}} \leftarrow \hat{\tilde{\mathbf{x}}} \cdot x_{\mathrm{scale}} - x_{\mathrm{off}}$
\State \Return $\hat{\mathbf{x}}$
\end{algorithmic}
\end{algorithm}

Default parameters are $\eta = 0.95$, $K_{\min} = 10$, $K_{\max} = 40$.
For narrowband modulations (BPSK) at high SNR, the selected $K$ is typically
10--15; for wideband digital modulations (QAM64, GFSK), it is 20--35.
On the RML2016.10a dataset, adaptive-$K$ achieves the highest overall
accuracy among all evaluated defenses, outperforming fixed Top-$K$ by 1--3\%
on CW attacks and 3--5\% on EAD attacks across SNR 0--18~dB.

\subsection{Classical Filter Baselines}
\label{sec:classical}

To contextualize the performance of our frequency-domain adaptive methods,
we evaluate five classical signal processing filters as adversarial defense
baselines. These represent the standard signal processing toolkit that a
communications engineer would apply before considering ML-specific defenses.
All filters are calibrated per-SNR via grid search over their hyperparameters
on a held-out validation subset of RML2016.10a, maximizing post-filter
classification accuracy on the undefended model.

\textbf{Kalman filter.} The Kalman filter~\cite{kalman1960new} models each IQ
channel as a first-order random walk corrupted by observation noise, with
process noise covariance $Q$ and observation noise covariance $R$ as tunable
parameters. The scalar filter recursively estimates the signal state using
prediction and update steps. For AMC defense, $R$ is calibrated to the noise
floor of adversarial perturbations. Kalman filtering is optimal for linear
Gaussian noise but was not designed for structured adversarial perturbations;
our experiments test whether its denoising properties generalize.

\textbf{Wiener filter.} The Wiener filter~\cite{wiener1949extrapolation}
computes the MMSE estimate in the frequency domain by applying a
signal-to-noise weighting $H_k = S_{ss,k} / (S_{ss,k} + S_{nn,k})$ to each
spectral bin, where $S_{ss,k}$ and $S_{nn,k}$ are the signal and noise power
spectral densities at frequency $k$, estimated from calibration data.
The Wiener filter assumes stationary signal and noise statistics; adversarial
perturbations violate the stationarity assumption.

\textbf{Savitzky-Golay filter.} The Savitzky-Golay filter~\cite{savitzky1964smoothing}
smooths each IQ channel by fitting a polynomial of degree $d$ to a sliding
window of $w$ samples using local least squares. The filter preserves the
shape of peaks and pulses better than standard lowpass filters, making it
suitable for preserving QAM constellation transitions. Parameters $d$ and
$w$ are calibrated via grid search.

\textbf{Gaussian filter.} A Gaussian kernel $h[n] = \exp(-n^2 / 2\sigma^2)$
is convolved with each IQ channel, implementing a smooth lowpass response with
cutoff determined by the standard deviation $\sigma$. The filter is implemented
as a GPU-native \texttt{F.conv1d} operation for efficient batch processing,
avoiding CPU roundtrip latency.

\textbf{FIR lowpass filter.} A finite impulse response (FIR) lowpass filter
with normalized cutoff frequency $f_c$ and filter order $M$ is applied to
each IQ channel~\cite{proakis2006digital}. The filter coefficients are computed
via windowed-sinc design (Hamming window). Like the Gaussian filter, the FIR
filter is implemented as a depthwise \texttt{F.conv1d} convolution for GPU
efficiency.

We also evaluate \textbf{randomized smoothing}~\cite{cohen2019certified} as a
certified robustness baseline. Randomized smoothing adds isotropic Gaussian
noise $\mathcal{N}(0, \sigma^2 I)$ to the input and takes the majority vote
over $k$ noisy copies. The certified radius is $\sigma \Phi^{-1}(p_A)$, where
$p_A$ is the probability of the top class under smoothing. We use $\sigma = 0.1$
and $k = 100$ copies, which is the lowest latency configuration that provides
non-trivial certificates. All five classical filters and randomized smoothing
are calibrated and evaluated under identical experimental conditions.

\subsection{Computational Complexity}
\label{sec:complexity}

The computational cost of each defense is dominated by its core operation.
FFT Top-$K$ (and adaptive-$K$) requires $\mathcal{O}(T \log T)$ for the real
FFT, $\mathcal{O}(T)$ for magnitude computation, $\mathcal{O}(T \log K)$ for
top-$K$ selection via partial sort, and $\mathcal{O}(T \log T)$ for the inverse
FFT. The spectral flatness computation adds $\mathcal{O}(T)$. Total per-sample
cost is $\mathcal{O}(T \log T)$, dominated by the FFT.

Among the classical filters, Gaussian and FIR convolutions cost
$\mathcal{O}(T \cdot M)$ per channel, where $M$ is the filter length (typically
17--33 taps); both are implemented as GPU-native depthwise convolutions with
negligible latency. Savitzky-Golay uses a precomputed coefficient vector and has
the same $\mathcal{O}(T \cdot w)$ complexity as a FIR filter. Kalman filtering
requires a sequential state update loop, with $\mathcal{O}(T)$ operations per
sample but poor parallelism (the state update at time $t$ depends on time $t-1$),
necessitating a CPU fallback and incurring higher per-batch latency despite lower
asymptotic complexity. Wiener filtering applies per-bin frequency-domain
weighting in $\mathcal{O}(T \log T + T)$ using a precomputed gain vector.
Randomized smoothing requires $k$ forward passes through the classifier ($k=100$
by default), making it $100\times$ more expensive than a single inference.
Detailed latency measurements and throughput comparisons are reported in
Section~\ref{sec:results}.
